\apendice{TODO: Documentación técnica de programación}

\section{Introducción}

Este manual se completará cuando exista una implementación estable. Documentarlo antes obligaría a describir comandos, módulos y dependencias que todavía pueden cambiar. El anexo conserva la estructura que deberá rellenarse durante la fase de entrega.

\section{Estructura de directorios}

Pendiente de reflejar la estructura real de los repositorios del frontend, backend, workers, infraestructura, pruebas y documentación.

\section{Configuración del entorno}

Se documentarán versiones compatibles de Java, Maven, Node.js, Docker y herramientas auxiliares, así como las variables necesarias. Los secretos no se incluirán en el repositorio; se proporcionarán ejemplos sin valores reales.

\section{Manual del programador}

Se describirán arquitectura, convenciones, módulos, contratos internos, migraciones, resolvers, colas, almacenamiento temporal y proceso para añadir una aplicación o un proveedor.

\section{Compilación, instalación y ejecución del proyecto}

Se incluirán procedimientos reproducibles para desarrollo local, pruebas, construcción de imágenes y despliegue. Los comandos se validarán desde una copia limpia del repositorio.

\section{Pruebas del sistema}

El manual final indicará cómo ejecutar pruebas unitarias, integración, extremo a extremo, seguridad y carga; qué servicios de prueba se requieren; y cómo interpretar los informes.

\section{Operación y resolución de incidencias}

Se documentarán métricas, logs, colas, limpieza de temporales, rotación de secretos, copias de seguridad, fuentes rotas y recuperación de trabajos. Esta sección se escribirá a partir del comportamiento observado, no de supuestos.
